% CS 418: Intro to Data Science — Worksheet 11 (11-week06-day1)
% Week 6, Monday 2026-09-28
% NOTICE: This worksheet was drafted by an LLM coding system from the
% instructor's list of checkpoints. It is maintained, reviewed, and vetted by humans.
% Compile with: latexmk -pdf worksheet.tex

\documentclass[11pt]{article}

% ── Packages ──────────────────────────────────────────────────────────────────
\usepackage[margin=0.65in, headheight=14pt]{geometry}
\usepackage{amsmath, amssymb}
\usepackage{ragged2e}
% Never break a word across lines, including at hyphens ("Sans-Serif").
\hyphenpenalty=10000
\exhyphenpenalty=10000
\usepackage{enumitem}
\usepackage{booktabs}
\usepackage{graphicx}
\usepackage{hyperref}
\usepackage{xcolor}
\usepackage{fancyhdr}
\usepackage{mdframed}
\usepackage{array}

% ── Draft watermark ───────────────────────────────────────────────────────────
% The watermark marks the PDF as unfinished so nobody mistakes a draft for the
% real thing. Delete this block once the worksheet has been reviewed.
\usepackage{draftwatermark}
\SetWatermarkText{IN-PROGRESS}
\SetWatermarkScale{2.2}
\SetWatermarkLightness{0.85}


% ── Page style ────────────────────────────────────────────────────────────────
\pagestyle{fancy}
\fancyhf{}
\lhead{\textbf{CS 418: Intro to Data Science}}
\rhead{UIC, Fall 2026}
\cfoot{\thepage}
\renewcommand{\headrulewidth}{1pt}
\setlength{\headheight}{12pt}
\setlength{\headsep}{0.4cm}

% ── Hyperlinks ────────────────────────────────────────────────────────────────
\hypersetup{colorlinks=false, hidelinks}

% ── Convenience environments ──────────────────────────────────────────────────
% Usage: \blankspace{6cm} — blank area for hand-written work, no rules
\newcommand{\blankspace}[1]{%
  \vspace{0.3em}
  \begin{minipage}[t][#1][t]{\linewidth}
  \end{minipage}
  \vspace{0.3em}
}

% Usage: \guidedopen{title}{guided content}{guided workspace height}{open-ended workspace height}
% Left: a titled, structured prompt with its own blank workspace. Right: a
% fully blank workspace for open-ended exploration of the same topic (no
% prompt text), sized so the two columns land at roughly the same total
% height. Neither workspace is ruled.
\newcommand{\guidedopen}[4]{%
  \noindent\textbf{#1}\\[0.18em]
  \noindent
  \begin{minipage}[t]{0.485\textwidth}
    #2
    \blankspace{#3}
  \end{minipage}
  \hfill
  \begin{minipage}[t]{0.485\textwidth}
    \blankspace{#4}
  \end{minipage}
  \\[0.4em]
}

% Usage: \panelbox{width}{height}{caption} — a framed drawing panel with `caption`
% centered beneath it. An empty caption leaves the panel bare, with nothing
% below the frame at all.
\newcommand{\panelbox}[3]{%
  \begin{minipage}[t]{\dimexpr#1+2\fboxsep+2\fboxrule\relax}
    \framebox[#1][c]{\rule{0pt}{#2}}%
    \if\relax\detokenize{#3}\relax\else
      \\[0.25em]{\centering\small #3\par}%
    \fi
  \end{minipage}%
}

% ══════════════════════════════════════════════════════════════════════════════
\begin{document}
\RaggedRight

% ── Title block: topic left, info box right ───────────────────────────────────
\noindent
\begin{minipage}[t]{0.55\textwidth}
  \vspace{0pt}
  {\LARGE \textbf{Worksheet \#\,11}} \\[0.18em]
  {\large \textit{Week 6, Day 1: Visualization rhetoric, uncertainty, and time series}}
\end{minipage}
\hfill
\begin{minipage}[t]{0.40\textwidth}
  \vspace{0pt}
  \small
  Name: \underline{\hspace{4.5cm}} \\[0.35em]
  NetID: \underline{\hspace{4.3cm}} \\[0.35em]
  Date: \underline{\hspace{4.4cm}} \\
\end{minipage}

\vspace{0.6em}
\noindent\rule{\linewidth}{0.6pt}
\vspace{0.1em}

\noindent\underline{\textbf{Guided checkpoints}}\\[0.2em]
\noindent{\small For full credit, fill in the checkpoints. If you skip a checkpoint, you must include your own notes for full credit.}
\vspace{0.8em}

% ── 1. Editorial layers ──────────────────────────────────────────────────────
\guidedopen{1. Four Editorial Layers of Rhetoric}{
List the four layers where a designer's choices can shape a visualization's message, with one example choice for each.
}{2.7cm}{3.6cm}

\vfill

% ── 2. Serif vs. sans-serif ──────────────────────────────────────────────────
\guidedopen{2. Serif vs.\ Sans-Serif}{
Write the word ``Data'' in each style. Circle one serif.

\vspace{0.45em}

\noindent
\panelbox{3.7cm}{1.6cm}{serif}\hfill
\panelbox{3.7cm}{1.6cm}{sans-serif}
}{0.2cm}{3.2cm}

\vfill

% ── 3. Visualizing intervals ─────────────────────────────────────────────────
\guidedopen{3. Four Ways to Visualize Uncertainty}{
Name the four kinds of displays for intervals and uncertainty, with one example of each.
}{2.7cm}{3.6cm}

\vfill

% ── 4. Reading the cone ──────────────────────────────────────────────────────
\guidedopen{4. Reading the Cone}{
Before the discussion: what do you think the cone represents? After: what does it actually represent?
}{3.0cm}{4.0cm}

\vfill

\newpage

% ── 5. Alternatives to the cone ──────────────────────────────────────────────
\guidedopen{5. Alternatives to the Cone}{
Name or sketch two other ways to show where a hurricane might go.
}{3.0cm}{3.9cm}

\vfill

% ── 6. Time series examples ──────────────────────────────────────────────────
\guidedopen{6. Common Examples of Time Series}{
List four examples of a time series: the same variable, measured over time.
}{3.0cm}{3.9cm}

\vfill

% ── 7. Four-part decomposition ───────────────────────────────────────────────
\guidedopen{7. Four-Part Decomposition}{
List the four parts of a time series, with a few words on what each one is.
}{3.0cm}{3.9cm}

\vfill

% ── 8. Additive vs. multiplicative ───────────────────────────────────────────
\guidedopen{8. Additive or Multiplicative?}{
Sketch one series with additive seasonality and one with multiplicative seasonality.

\vspace{0.45em}

\noindent
\panelbox{7.8cm}{2.2cm}{additive}

\vspace{0.5em}

\noindent
\panelbox{7.8cm}{2.2cm}{multiplicative}
}{0.2cm}{6.6cm}

\vfill

\end{document}
